\chapter{Objetivo do manual}
Este manual foi escrito para membros e utilizadores operacionais da Smarter Hub. O objetivo é que qualquer pessoa consiga usar a plataforma sem depender de alguém da equipa técnica para cada ação simples.

A Smarter Hub reúne informação de perfil, pedidos, férias, ausências, formações, equipa, notificações e saúde e bem-estar. Em termos práticos, o portal substitui vários pedidos dispersos por um só ponto de entrada.

\chapter{Como ler este manual}
\section{Estrutura}
Os capítulos estão organizados por ecrã, como aparecem no portal. Sempre que houver um fluxo completo, ele é explicado passo a passo. Quando o nome do ecrã muda entre PT-PT e PT-BR, o texto apresenta os dois nomes.

\section{Notas de nomenclatura PT-PT e PT-BR}
\begin{longtable}{|p{0.34\textwidth}|p{0.58\textwidth}|}
\hline
\textbf{PT-PT} & \textbf{PT-BR / observação} \\
\hline
A Minha Ficha & Em PT-BR, costuma ser entendido como \textit{Meu cadastro} ou \textit{Meu perfil}. Na plataforma a designação visível é A Minha Ficha. \\
\hline
Formações & Em PT-BR também pode aparecer como \textit{Treinamentos} no discurso interno. \\
\hline
Chefes de equipa & Em PT-BR, pode surgir como \textit{líderes de equipe} ou \textit{gestores de equipe}. \\
\hline
Férias / Ausências & O nome é mantido, mas os tipos de ausência e a linguagem de apoio podem variar por país. \\
\hline
Saúde e bem-estar & Nome mantido, com conteúdos e documentação ajustados por país. \\
\hline
Membros & Em PT-BR, o termo é o mesmo, mas algumas descrições podem usar \textit{pessoas} ou \textit{usuários internos}. \\
\hline
\end{longtable}

\chapter{Percurso do membro}
\section{Entrar na plataforma}
Ao entrar, o sistema carrega a conta autenticada e apresenta a página inicial. Se a empresa usar integração com Microsoft, o acesso é feito com a conta autorizada pela organização.

\section{O que o membro vê primeiro}
\begin{itemize}
  \item Atalhos para as áreas mais usadas.
  \item Pendências ou avisos relevantes.
  \item Notificações recentes.
  \item O menu lateral com as secções a que tem acesso.
\end{itemize}

\section{Regras de uso básicas}
\begin{itemize}
  \item Confirme sempre se está no perfil certo antes de abrir um pedido.
  \item Guarde alterações antes de fechar o separador.
  \item Leia as notificações até ao fim; elas explicam o que já foi aprovado ou rejeitado.
  \item Se um campo não estiver editável, isso pode ser uma regra de aprovação e não um erro.
\end{itemize}

\chapter{Página inicial}
\section{Home}
A Home é o ponto de entrada. A copy pode mudar conforme o perfil, mas a lógica é sempre a mesma: dar visibilidade rápida ao que importa naquele momento.

\subsection{Elementos normalmente visíveis}
\begin{itemize}
  \item Cartões de pendências.
  \item Resumo de formações ativas.
  \item Notificações.
  \item Atalhos para ficha, equipas, férias, formações e bem-estar.
\end{itemize}

\subsection{Fluxo de utilização}
\begin{enumerate}
  \item Entre na Home.
  \item Identifique os cartões mais relevantes.
  \item Abra o cartão ou atalho que corresponde ao trabalho que precisa de fazer.
  \item Use a Home como ponto de retorno quando estiver a navegar por outras áreas.
\end{enumerate}

\chapter{Ficha pessoal}
\section{A Minha Ficha / Meu cadastro}
Esta é a sua área pessoal. É aqui que consulta e, quando permitido, atualiza informação que a empresa usa para suporte interno, contacto e validação administrativa.

\subsection{O que costuma existir}
\begin{itemize}
  \item Dados pessoais.
  \item Contactos.
  \item Documentos de identificação.
  \item Dados fiscais e bancários.
  \item Contacto de emergência.
  \item Dados contratuais.
  \item Dados de formação.
  \item Pedido de benefícios.
\end{itemize}

\subsection{O que cada bloco significa}
\begin{itemize}
  \item \textbf{Dados pessoais}: nome, data de nascimento, nacionalidade e outros dados-base.
  \item \textbf{Dados de contacto}: email, telefone e morada.
  \item \textbf{Documentos}: identificações, números oficiais e validade.
  \item \textbf{Dados fiscais e bancários}: informação usada para processos administrativos.
  \item \textbf{Contacto de emergência}: pessoa a contactar em caso de necessidade.
  \item \textbf{Dados contratuais}: contexto laboral e informações associadas ao vínculo.
  \item \textbf{Dados de formação}: histórico e qualificações.
  \item \textbf{Benefícios}: pedidos ou dados associados a benefícios internos.
\end{itemize}

\subsection{Como editar corretamente}
\begin{enumerate}
  \item Abra \textit{A Minha Ficha}.
  \item Vá à secção correta.
  \item Corrija apenas o campo pretendido.
  \item Leia a mensagem do sistema antes de gravar.
  \item Guarde as alterações e confirme a notificação de sucesso.
\end{enumerate}

\subsection{O que pode acontecer depois de guardar}
Dependendo do tipo de campo, a alteração pode ficar logo refletida ou seguir para validação/aprovação. Quando isso acontece, a sua alteração deixa de ser apenas uma edição local e passa a fazer parte de um fluxo formal.

\chapter{Percurso e desenvolvimento}
\section{Plano de Carreira}
O Plano de Carreira ajuda a perceber como a empresa organiza a progressão, as expectativas e as etapas de crescimento profissional.

\subsection{Ecrãs e conteúdo}
\begin{itemize}
  \item \textbf{O meu nível}: contexto atual e posição no percurso.
  \item \textbf{Roadmap de níveis}: visão de etapas futuras.
  \item \textbf{Processo de avaliação}: ligação ao ciclo de avaliação ou checkpoints.
\end{itemize}

\subsection{Como usar}
\begin{enumerate}
  \item Abra o Plano de Carreira.
  \item Leia o nível atual e o próximo objetivo.
  \item Abra o roadmap para perceber o caminho de progressão.
  \item Consulte o processo de avaliação para entender o que é esperado em cada fase.
\end{enumerate}

\chapter{Trabalho em equipa}
\section{Equipas}
A página de Equipas mostra a estrutura organizacional a que pertence. Serve para perceber quem faz parte da mesma equipa, quem lidera e como a equipa está distribuída.

\subsection{O que pode ver}
\begin{itemize}
  \item Nome da equipa.
  \item Pessoas associadas.
  \item Relações hierárquicas quando existem.
  \item Dados úteis para planeamento de trabalho.
\end{itemize}

\subsection{Situações típicas de uso}
\begin{itemize}
  \item Confirmar quem está na sua equipa.
  \item Ver quem pode responder por um pedido.
  \item Identificar rapidamente a pessoa certa para um assunto operacional.
\end{itemize}

\chapter{Formação}
\section{Formações}
Este módulo junta o histórico e o estado das formações. Dependendo das permissões, pode mostrar apenas as suas formações ou também as da equipa.

\subsection{O que a página costuma mostrar}
\begin{itemize}
  \item Lista de formações.
  \item Estado de cada formação.
  \item Data de início.
  \item Entidade que atribuiu ou registou a formação.
  \item Filtros por estado, origem, entidade, horas e datas.
\end{itemize}

\subsection{Fluxos mais comuns}
\begin{enumerate}
  \item Abrir a página de Formações.
  \item Procurar a formação pelo nome ou pelo membro.
  \item Ler o detalhe para entender o estado.
  \item Marcar como concluída, se o fluxo permitir.
  \item Exportar para Excel, quando essa opção estiver disponível.
\end{enumerate}

\subsection{Leitura prática dos estados}
\begin{itemize}
  \item \textbf{ASSIGNED}: formação atribuída e ainda em curso.
  \item \textbf{COMPLETED}: formação concluída.
  \item Outros estados podem existir consoante a evolução da plataforma.
\end{itemize}

\chapter{Férias e ausências}
\section{Férias / Ausências}
Esta é uma das áreas mais usadas por qualquer membro. Serve para pedir férias, registar ausências e acompanhar o resultado do pedido.

\subsection{Separadores existentes}
\begin{itemize}
  \item \textbf{Resumo}: visão geral e saldos.
  \item \textbf{Calendário}: leitura por dias, semanas ou ano.
  \item \textbf{Férias da equipa}: leitura coletiva para evitar conflitos.
  \item \textbf{Dias automáticos}: manutenção dos dias dados pela empresa.
  \item \textbf{Mapa de Férias}: apoio à exportação e consulta consolidada.
  \item \textbf{Atribuição imediata}: usos rápidos quando o perfil o permite.
\end{itemize}

\subsection{Como pedir férias}
\begin{enumerate}
  \item Abra \textit{Férias / Ausências}.
  \item Entre no separador que corresponde ao fluxo pretendido.
  \item Escolha \textit{Férias} ou \textit{Ausência}, se o sistema lhe mostrar essa opção.
  \item Indique a data de início e a data de fim.
  \item Se for meio-dia, confirme o período parcial quando existir a opção.
  \item Adicione uma observação clara.
  \item Submeta o pedido.
\end{enumerate}

\subsection{Como acompanhar o pedido}
\begin{itemize}
  \item Abra o Resumo para ver o ponto de situação.
  \item Use o Calendário para perceber a posição temporal do pedido.
  \item Consulte as Notificações para ver aprovações ou rejeições.
  \item Leia a justificação quando houver uma rejeição ou pedido de correção.
\end{itemize}

\subsection{Dias automáticos}
Os dias automáticos representam dias dados pela empresa, por país e por ano. Normalmente não são criados pelo membro comum, mas aparecem no calendário para ajudar a leitura do período.

\section{Leitura das ausências}
\begin{itemize}
  \item Férias aprovadas aparecem no calendário.
  \item Ausências médicas e outras ausências podem ter regras próprias.
  \item O planeamento da equipa ajuda a evitar pedidos em simultâneo que criem impacto operacional.
\end{itemize}

\chapter{Notificações}
\section{Notificações}
As notificações funcionam como histórico curto e como ponto de entrada para decisões ou detalhes.

\subsection{Tipos de avisos comuns}
\begin{itemize}
  \item Pedido aprovado.
  \item Pedido rejeitado.
  \item Pedido a aguardar análise.
  \item Alteração de ficha.
  \item Informação operacional ligada a um módulo.
\end{itemize}

\subsection{Como usar}
\begin{enumerate}
  \item Clique no sino do topo.
  \item Leia a mensagem completa.
  \item Abra a notificação, se ela levar para outra área.
  \item Marque como lida quando já terminou o tratamento.
\end{enumerate}

\chapter{Saúde e bem-estar}
\section{Saúde e bem-estar}
Esta área foi desenhada para disponibilizar conteúdos de apoio e, em alguns casos, canais de reporte sensível.

\subsection{Conteúdos comuns nas duas geografias}
\begin{itemize}
  \item Formulário de assédio.
  \item Ergonomia.
  \item Suporte Básico de Vida.
\end{itemize}

\subsection{Diferença entre conteúdo comum e conteúdo local}
Alguns blocos aparecem em ambas as geografias porque são transversais. Outros podem existir apenas num país. Isto é importante para evitar confundir conteúdos gerais com conteúdos legal ou operacionalmente específicos.

\subsection{Como usar a página}
\begin{enumerate}
  \item Abra \textit{Saúde e bem-estar}.
  \item Identifique se o bloco que quer é PDF ou formulário.
  \item Abra o ficheiro ou o formulário.
  \item Se for um reporte sensível, escreva de forma objetiva e completa.
\end{enumerate}

\subsection{Quando usar o formulário de assédio}
Use este formulário quando precisar de reportar uma situação que exija acompanhamento formal. O texto deve ser claro, factual e suficientemente detalhado para permitir triagem.

\chapter{Chefes de equipa}
\section{O que muda quando é líder de equipa}
Chefes de equipa têm mais visibilidade sobre a equipa e, em alguns casos, mais capacidade de leitura e decisão. O manual não deve ser lido como uma versão diferente da plataforma, mas como uma camada adicional de responsabilidade.

\section{O que olhar primeiro}
\begin{itemize}
  \item Pedidos pendentes da equipa.
  \item Sobreposições de férias.
  \item Formações relevantes para a operação.
  \item Notificações que exigem resposta rápida.
\end{itemize}

\section{Como usar Equipas de forma eficaz}
\begin{enumerate}
  \item Abra a página de Equipas.
  \item Selecione a equipa certa.
  \item Verifique os membros e respetivas associações.
  \item Use a informação da página para tomar decisões mais informadas.
\end{enumerate}

\section{Como gerir férias da equipa}
\begin{enumerate}
  \item Entre em \textit{Férias / Ausências}.
  \item Abra o separador de equipa.
  \item Analise o calendário coletivo.
  \item Verifique feriados, dias automáticos e ausências já planeadas.
  \item Tome decisão com base no impacto operacional.
\end{enumerate}

\section{Como acompanhar formações da equipa}
\begin{itemize}
  \item Identifique formações em aberto.
  \item Veja quem atribuiu a formação.
  \item Confirme o estado de conclusão.
  \item Use a informação para alinhar formação com prioridades reais da equipa.
\end{itemize}

\chapter{Perguntas frequentes}
\section{Não vejo um módulo no menu}
Isso significa normalmente que o seu perfil não tem acesso a essa funcionalidade. A plataforma mostra apenas o que faz sentido para o seu contexto.

\section{O sistema não deixa alterar um campo}
Pode existir validação, bloqueio por estado do processo ou necessidade de aprovação.

\section{Não sei se usei o ecrã certo}
Compare o nome do menu com a nomenclatura deste manual. Se estiver entre A Minha Ficha, perfil e conta, lembre-se que são coisas diferentes.

\chapter{Anexo de ecrãs detalhados}
\section{Home}
Este é o ecrã de entrada. O objetivo não é editar nada, mas orientar o utilizador para a próxima ação.

\subsection{O que observar}
\begin{itemize}
  \item Mensagem de contexto no topo.
  \item Atalhos mais importantes do seu perfil.
  \item Cartões de estado para pendências e itens a acompanhar.
  \item Notificações recentes.
\end{itemize}

\subsection{Como usar bem a Home}
\begin{enumerate}
  \item Comece por ler os cartões que tiverem contagem maior.
  \item Abra primeiro o item que mais impacta o seu trabalho.
  \item Use a Home para regressar a um ponto neutro entre módulos.
\end{enumerate}

\section{A Minha Ficha}
Este ecrã é o seu registo pessoal. O valor dele está em manter os seus dados consistentes para o resto do processo funcionar sem fricção.

\subsection{Antes de alterar}
\begin{itemize}
  \item Confirme se o campo está mesmo errado.
  \item Veja se o campo é apenas informativo ou editável.
  \item Leia se o sistema indica pedido pendente, validação ou bloqueio.
\end{itemize}

\subsection{Após alterar}
\begin{itemize}
  \item Verifique a mensagem de confirmação.
  \item Confirme se o valor gravado ficou visível.
  \item Aguarde eventual aprovação, se o campo estiver sujeito a validação.
\end{itemize}

\section{Plano de Carreira}
\subsection{Leitura correta do conteúdo}
O Plano de Carreira é um ecrã de alinhamento: mostra estado atual, referenciais de evolução e o que a empresa espera em cada etapa.

\subsection{Boa prática}
\begin{itemize}
  \item Leia primeiro o nível atual.
  \item Consulte depois o roadmap.
  \item Use o processo de avaliação para perceber o que está associado ao momento presente.
\end{itemize}

\section{Equipas}
\subsection{Quando este ecrã é útil}
\begin{itemize}
  \item Quando quer perceber quem faz parte da mesma equipa.
  \item Quando quer saber quem lidera ou coordena.
  \item Quando precisa de contexto para férias, formações ou avisos.
\end{itemize}

\section{Formações}
\subsection{Campos e leitura}
Mesmo sem mapear cada coluna visual, o utilizador deve ler sempre quatro coisas: nome da formação, estado, data de início e origem. Essas quatro variáveis dizem quase tudo sobre o contexto do registo.

\subsection{Cenários comuns}
\begin{itemize}
  \item Formação própria criada pelo utilizador.
  \item Formação atribuída por alguém da empresa.
  \item Formação concluída.
  \item Formação ainda em curso.
\end{itemize}

\section{Férias / Ausências}
\subsection{Resumo}
O resumo deve ser lido como fotografia do momento. É o lugar onde se percebe saldo, situação atual e visão geral do ano.

\subsection{Calendário}
\begin{itemize}
  \item Use o calendário para evitar pedidos em cima de feriados ou dias já ocupados.
  \item Repare em dias automáticos, porque eles também contam na leitura do período.
  \item Compare o seu pedido com os da equipa quando existir esse separador.
\end{itemize}

\subsection{Dias automáticos}
Este ecrã é informativo para o membro, mas é operacionalmente sensível para RH. O membro deve entender que um dia automático não é um pedido individual, é uma regra da empresa aplicada ao calendário.

\section{Saúde e bem-estar}
\subsection{Leitura por tipo de bloco}
\begin{itemize}
  \item Se o bloco for PDF, a ação é abrir e ler.
  \item Se o bloco for formulário, a ação é preencher e submeter.
\end{itemize}

\subsection{Conteúdo transversal}
Formulário de assédio, ergonomia e suporte básico de vida devem ser tratados como conteúdos comuns nas duas geografias. A lógica do utilizador é a mesma em Portugal e no Brasil, mesmo que a apresentação ou o conteúdo documental possa variar.

\chapter{Guia de nomenclatura prática}
\section{Quando o nome muda, o significado pode não mudar}
\begin{longtable}{|p{0.3\textwidth}|p{0.6\textwidth}|}
\hline
\textbf{Nome} & \textbf{Significado prático} \\
\hline
A Minha Ficha & Ficha pessoal editável do membro. \\
\hline
Perfil / Conta & Identidade da conta Microsoft ou detalhe de autenticação. \\
\hline
Férias / Ausências & Área de pedidos, calendário e leitura de ausências. \\
\hline
Saúde e bem-estar & Biblioteca de apoio e reporte sensível. \\
\hline
Equipas & Estrutura de colaboração, membros e chefias. \\
\hline
Formações / Treinamentos & Histórico, atribuição e conclusão de formações. \\
\hline
\end{longtable}

\chapter{Conclusão}
Use este manual como referência de operação diária. O ideal é revisá-lo sempre que houver mudanças de processo, novos campos, novos separadores ou alterações de nomenclatura entre PT-PT e PT-BR.

\chapter{Anexo operacional por ecrã}
\section{Mapa rápido dos ecrãs principais}
\begin{longtable}{|p{0.18\textwidth}|p{0.14\textwidth}|p{0.17\textwidth}|p{0.21\textwidth}|p{0.24\textwidth}|}
\hline
\textbf{Ecrã} & \textbf{Quem usa} & \textbf{Abertura} & \textbf{Botões / abas-chave} & \textbf{Notas PT-PT / PT-BR} \\
\hline
Home & Todos & Rota inicial & Atalhos, cartões, notificações & A copy muda por perfil; no BR pode haver linguagem mais descritiva em avisos. \\
\hline
A Minha Ficha & Membro autorizado & \texttt{/profile} & Etapas, checklist, guardar, submeter pedido & PT-PT usa linguagem de ficha; PT-BR entende-se como cadastro/perfil. \\
\hline
Plano de Carreira & Membro autorizado & \texttt{/plano-carreira} & O meu nível, roadmap, avaliação & Conteúdo de progressão deve ser lido com a terminologia interna do país. \\
\hline
Equipas & Todos exceto CONVIDADO & \texttt{/equipas} & Modal de equipa, visão de membros, comparação & Chefias tendem a ver mais contexto; em PT-BR a leitura de equipe pode ser mais comum. \\
\hline
Formações & Todos os autorizados & \texttt{/formacoes} & Exportar Excel, Nova formação, filtros, concluir & Em PT-BR o utilizador pode chamar o módulo de treinamentos; a UI mantém Formações. \\
\hline
Férias / Ausências & Todos os autorizados & \texttt{/ferias} & Resumo, Calendário, Férias da equipa, Dias automáticos, Mapa, Atribuição imediata & Há diferenças por país nos motivos e na interpretação de dias automáticos. \\
\hline
Saúde e bem-estar & Todos exceto CONVIDADO & \texttt{/saude-bem-estar} & Editar página, Reportar situação, PDFs, formulários & Formulário de assédio, ergonomia e suporte básico de vida são comuns a PT e BR. \\
\hline
Notificações & Todos & \texttt{/notifications} & Filtrar, abrir, marcar como lida, apagar & A barra superior concentra o acesso; a língua da mensagem pode ajustar termos técnicos. \\
\hline
\end{longtable}

\section{Home}
\subsection{Objetivo real}
A Home não é uma página decorativa. É a leitura do estado operacional do utilizador no momento em que entra na plataforma.

\subsection{O que aparece}
\begin{itemize}
  \item Cartões de pendências.
  \item Entradas rápidas para módulo críticos.
  \item Notificações recentes.
  \item Contexto resumido do perfil atual.
\end{itemize}

\subsection{Exemplo de uso}
Se o cartão de férias pendentes estiver com contagem elevada, o utilizador deve abrir Férias / Ausências antes de continuar a explorar outras áreas.

\subsection{Erro típico}
Entrar na Home e ignorar os cartões de pendência. Isso costuma atrasar a resolução de pedidos já em curso.

\section{A Minha Ficha}
\subsection{Fluxo de ponta a ponta}
\begin{enumerate}
  \item Abrir a ficha pessoal.
  \item Navegar entre as secções.
  \item Corrigir apenas o que estiver necessário.
  \item Guardar.
  \item Aguardar validação, se aplicável.
\end{enumerate}

\subsection{O que muda entre PT-PT e PT-BR}
\begin{itemize}
  \item Em PT-PT é comum falar em NIF, NISS, cartão de cidadão e IBAN.
  \item Em PT-BR aparecem CPF, PIS, CTPS, RG e CNH.
  \item A estrutura é a mesma, mas o conteúdo fiscal e documental muda consoante o país.
\end{itemize}

\subsection{Exemplo de uso}
Um membro muda de morada. Abre A Minha Ficha, entra em dados de contacto, altera a morada habitual/fiscal e confirma a gravação.

\subsection{Erro típico}
Editar uma morada por cima da errada quando o mesmo utilizador tem campos distintos para país ou contexto fiscal.

\section{Plano de Carreira}
\subsection{Como ler corretamente}
O plano de carreira deve ser lido como documento de orientação, não como lista de tarefas imediatas.

\subsection{Exemplo de uso}
Um membro quer perceber o que significa o seu nível atual. Abre o separador correspondente e compara o roadmap com o processo de avaliação.

\subsection{Erro típico}
Assumir que o roadmap substitui a conversa com a chefia. Ele complementa, não substitui.

\section{Equipas}
\subsection{Como usar o detalhe de equipa}
\begin{itemize}
  \item Confirmar quem pertence à equipa.
  \item Observar lideranças.
  \item Usar o modal de detalhe quando existir uma análise mais profunda.
\end{itemize}

\subsection{Diferença PT-PT / PT-BR}
Em Portugal é comum falar em \textit{equipa}; no Brasil, \textit{equipe}. A plataforma usa equipa como referência visual, mas o manual pode adaptar a leitura conforme a audiência.

\subsection{Erro típico}
Interpretar a equipa como uma lista estática. Na prática, ela deve ser usada como contexto vivo para férias, formações e decisões.

\section{Formações}
\subsection{Campos e botões que importam}
\begin{itemize}
  \item \textbf{Nova formação}: cria um novo registo.
  \item \textbf{Exportar Excel}: extrai a informação para análise.
  \item \textbf{Concluir}: fecha o ciclo de uma formação.
  \item \textbf{Filtros}: reduzem ruído e ajudam a encontrar o registo certo.
\end{itemize}

\subsection{Leitura por país}
\begin{itemize}
  \item O nome visível é Formações, mas em PT-BR muitos utilizadores pensarão em Treinamentos.
  \item O conceito funcional é o mesmo: atribuir, acompanhar, concluir.
\end{itemize}

\subsection{Exemplo de uso}
O membro procura todas as formações ativas do último ano. Filtra por estado e lê a coluna de origem para entender quem criou o registo.

\subsection{Erro típico}
Confundir uma formação própria com uma atribuída e tentar corrigir um estado que não deveria ser alterado manualmente.

\section{Férias / Ausências}
\subsection{Fluxos por separador}
\begin{itemize}
  \item \textbf{Resumo}: fotografia do saldo e da situação corrente.
  \item \textbf{Calendário}: leitura de dias e sobreposições.
  \item \textbf{Férias da equipa}: visão coletiva.
  \item \textbf{Dias automáticos}: política anual da empresa.
  \item \textbf{Mapa de Férias}: consolidação.
  \item \textbf{Atribuição imediata}: registo mais rápido.
\end{itemize}

\subsection{Diferenças país a país}
\begin{itemize}
  \item Em Portugal é mais comum a leitura de feriados locais, tolerâncias e políticas ligadas ao contexto PT.
  \item No Brasil surgem escolhas de âmbito que podem incluir estados específicos, como SP ou RS.
  \item Os tipos de ausência e a linguagem usada nos motivos podem variar entre países.
\end{itemize}

\subsection{Exemplo de uso}
O membro quer férias em agosto. Começa pelo resumo para ver saldo, depois consulta o calendário e, por fim, submete o pedido com observação curta.

\subsection{Erro típico}
Pedir datas sem cruzar com o calendário da equipa. Isso aumenta a probabilidade de rejeição ou de conflito operacional.

\section{Saúde e bem-estar}
\subsection{Leitura de blocos}
\begin{itemize}
  \item PDF: ler e abrir documento.
  \item Formulário: preencher e submeter.
\end{itemize}

\subsection{O que é transversal}
Formulário de assédio, ergonomia e suporte básico de vida são conteúdos que devem aparecer iguais em Portugal e no Brasil.

\subsection{Exemplo de uso}
O membro procura o formulário de reporte de assédio, lê as instruções e submete uma descrição factual com contacto preferencial, se desejar.

\subsection{Erro típico}
Tratar conteúdos sensíveis como textos informativos normais. Alguns blocos exigem mais cuidado, mais contexto e mais formalidade.

\section{Notificações}
\subsection{Como interpretar}
As notificações devem ser lidas como uma lista curta de ações ou mudanças relevantes. Se a mensagem for técnica, use o contexto do título e do botão para perceber o que aconteceu.

\subsection{Exemplo de uso}
Uma notificação avisa que uma alteração de ficha foi rejeitada. O membro abre a notificação, lê o motivo e regressa à ficha para corrigir o campo indicado.

\subsection{Erro típico}
Ignorar notificações e tentar descobrir o estado do pedido noutro módulo, sem abrir a mensagem original.

\chapter{Exemplos e erros frequentes}
\section{Exemplos práticos de utilização}
\begin{itemize}
  \item Um membro quer saber se o pedido de férias foi aprovado: consulta Notificações e depois Calendário.
  \item Um membro quer atualizar o telefone: abre A Minha Ficha, altera o campo e guarda.
  \item Um membro quer consultar conteúdos de ergonomia: entra em Saúde e bem-estar e abre o bloco correspondente.
  \item Um membro quer ver o seu percurso: entra em Plano de Carreira e lê o roadmap.
\end{itemize}

\section{Erros mais comuns}
\begin{itemize}
  \item Confundir conta Microsoft com ficha pessoal.
  \item Ignorar o calendário antes de pedir férias.
  \item Procurar um documento local quando o bloco é comum às duas geografias.
  \item Ler uma notificação sem abrir a página associada.
\end{itemize}

\section{Checklist mental antes de agir}
\begin{enumerate}
  \item Estou no ecrã certo?
  \item O meu país/perfil é o correto?
  \item O bloco ou pedido é comum ou específico do meu país?
  \item Preciso de guardar, submeter, aprovar ou apenas consultar?
\end{enumerate}

\chapter{Edição Portugal: operação local detalhada}
\section{Campos e documentação mais relevantes em PT}
Em Portugal, o membro deve prestar atenção especial aos campos fiscais e de identificação seguintes:
\begin{itemize}
  \item NIF.
  \item NISS.
  \item Cartão de Cidadão e validade.
  \item Registo criminal (ficheiro em PDF/JPG).
  \item IBAN e comprovativo associado.
  \item Situação IRS e declaração IRS.
\end{itemize}

\begin{manualnote}[Quando atualizar a ficha em PT]
Atualize primeiro os campos de identificação, depois os fiscais e por fim os comprovativos. Este ordenamento reduz erros de validação no fluxo de aprovação.
\end{manualnote}

\section{Férias em Portugal}
\subsection{Como pensar o pedido}
\begin{enumerate}
  \item Confirmar saldo no Resumo.
  \item Verificar feriados e dias automáticos no Calendário.
  \item Cruzar com Férias da equipa para evitar sobreposição crítica.
  \item Submeter pedido com observação objetiva.
\end{enumerate}

\subsection{Erros frequentes em PT}
\begin{itemize}
  \item Submeter férias sem confirmar impacto na equipa.
  \item Ignorar dias automáticos já registados pela empresa.
  \item Pedir período com datas incoerentes (início/fim).
\end{itemize}

\section{Saúde e bem-estar em Portugal}
\subsection{Conteúdos transversais obrigatórios}
\begin{itemize}
  \item Formulário de assédio.
  \item Ergonomia.
  \item Suporte Básico de Vida.
\end{itemize}

\subsection{Conteúdos locais possíveis}
\begin{itemize}
  \item Política interna de apoio social em PT.
  \item Procedimentos locais de contacto e apoio.
  \item Documentos de cobertura de saúde locais.
\end{itemize}

\chapter{Portugal: fluxos completos exemplificados}
\section{Fluxo A: corrigir IBAN e comprovativo}
\begin{enumerate}
  \item Abrir A Minha Ficha.
  \item Ir a Dados Fiscais e Bancários.
  \item Atualizar IBAN.
  \item Carregar comprovativo IBAN.
  \item Guardar e confirmar a mensagem de sucesso.
  \item Verificar notificação caso exista aprovação pendente.
\end{enumerate}

\section{Fluxo B: pedido de férias anual}
\begin{enumerate}
  \item Abrir Férias / Ausências.
  \item Entrar em Resumo e validar saldo.
  \item Abrir Calendário e analisar datas.
  \item Criar pedido com período completo.
  \item Submeter.
  \item Acompanhar em Notificações até decisão final.
\end{enumerate}

\section{Fluxo C: reporte sensível}
\begin{enumerate}
  \item Abrir Saúde e bem-estar.
  \item Selecionar Formulário de assédio.
  \item Descrever o caso de forma factual.
  \item Incluir contacto preferencial quando necessário.
  \item Submeter e aguardar retorno via fluxo interno.
\end{enumerate}
